Medical image segmentation is a critical component in assisting disease diagnosis and treatment planning. However, the development of deep learning methods in this field has long been constrained by the scarcity of high-quality annotated data. The clinical annotation process is not only time-consuming and labor-intensive, but the distribution of data also varies significantly across different medical institutions. This leads to poor generalization performance of existing models in cross-center and cross-modal scenarios, making them difficult to deploy effectively in real-world applications.

To address these two core challenges—annotation scarcity and insufficient model generalization—this study systematically constructs an adaptive theoretical and methodological framework for medical image segmentation. On one hand, to improve the utilization efficiency of limited annotated data, a cognitive uncertainty-guided semi-supervised learning framework is designed, which optimizes pseudo-label quality by quantifying the uncertainty of model predictions. On the other hand, to enhance the model's ability to extract multi-scale features, the fusion paradigm of U-Net and Transformer is thoroughly investigated, and a channel attention mechanism is introduced into the skip connections of the network to reinforce the transmission of critical features.

To validate the effectiveness of the proposed methods, a series of comparative experiments are conducted on two datasets: the CA dataset and the Synapse dataset. In the semi-supervised experiment, the proposed UA-MT framework achieves a Dice coefficient of 88.18\% using only 20\% of labeled data (16 cases), significantly outperforming the standard Mean Teacher method. In the network architecture experiment, the DA-TransUNet model, which incorporates a dual attention mechanism, attains an average Dice coefficient of 80.67\% on the multi-organ segmentation task, with boundary segmentation accuracy (HD95) clearly surpassing baseline models such as U-Net and Swin-Unet. These experimental results demonstrate that the proposed methods can effectively alleviate the bottleneck caused by scarce annotated data while substantially improving segmentation accuracy and model robustness.